% ====================================================================
% FF2-1 (ch2 §2.8 배치 제안) — 실계산 ∂U_oc/∂T(x̄) 완전식·단순식 + 가역 발열 부호 교대
%   좌표 = Chapter 1 staging 4-전이 파라미터(ΔH_rxn/ΔS_rxn/Q, w=RT/F, 298.15 K)로
%   전하 보존 음함수 Σ Q_j ξ_j(U_oc)=Q x̄ 를 풀어 eq:complete / eq:weighted / eq:qrev 를
%   점별 수치 평가(python)한 값. anchor: tab:qrev 5점·계산 예제(-0.204/-0.134 mV/K) 재현.
%   사용 라이브러리: ch2 preamble 범위 내(arrows.meta 만 사용).
% ====================================================================
\begin{figure}[t]
\centering
\begin{tikzpicture}
% =================== (a) dU_oc/dT(xbar) ===================
\begin{scope}[x=6.9cm,y=4.6cm]
% axes
\draw[-{Latex[length=1.6mm]}] (-0.02,-0.47) -- (-0.02,0.44)
  node[above,font=\scriptsize] {$\partial U_\oc/\partial T$ [mV/K]};
\draw[-{Latex[length=1.6mm]}] (-0.02,-0.47) -- (1.06,-0.47)
  node[below,font=\scriptsize] {$\bar x$};
\foreach \xx in {0.2,0.4,0.6,0.8,1.0}
  {\draw (\xx,-0.463) -- (\xx,-0.477) node[below,font=\scriptsize] {\xx};}
\foreach \yy/\lab in {-0.4/-0.4,-0.2/-0.2,0/0,0.2/0.2,0.4/0.4}
  {\draw (-0.013,\yy) -- (-0.027,\yy) node[left,font=\scriptsize] {\lab};}
% zero line
\draw[densely dashed,black!45] (-0.02,0) -- (1.04,0);
% transition levels dS0_j/F (dominance windows)
\draw[densely dotted,thick] (0.00,-0.166) -- (0.40,-0.166);
\node[font=\tiny,anchor=west] at (0.01,-0.196) {$\Delta S^0_{2\to1}/F=-0.166$};
\draw[densely dotted,thick] (0.53,-0.052) -- (0.75,-0.052);
\node[font=\tiny,anchor=west] at (0.40,-0.075) {$-0.052$ ($2\mathrm L{\to}2$)};
\node[font=\tiny,anchor=west] at (0.755,-0.025) {$0$ ($3{\to}2\mathrm L$)};
\draw[densely dotted,thick] (0.88,0.301) -- (1.03,0.301);
\node[font=\tiny,anchor=east] at (0.88,0.301) {$+0.301$ ($4{\to}3$)};
% sign-flip marker
\draw[densely dotted,black!55] (0.681,-0.47) -- (0.681,0.10);
\node[font=\tiny,black!55,anchor=south,rotate=90] at (0.681,-0.40) {부호 반전 $\bar x\!\approx\!0.68$};
% simple formula (centers only), dashed black
\draw[thick,densely dashed] plot[smooth] coordinates {(0.050,-0.146) (0.075,-0.145) (0.100,-0.144) (0.125,-0.143) (0.150,-0.141) (0.175,-0.139) (0.200,-0.138) (0.225,-0.136) (0.250,-0.134) (0.275,-0.132) (0.300,-0.130) (0.325,-0.127) (0.350,-0.125) (0.375,-0.122) (0.400,-0.119) (0.425,-0.116) (0.450,-0.112) (0.475,-0.109) (0.500,-0.105) (0.525,-0.101) (0.550,-0.097) (0.575,-0.093) (0.600,-0.089) (0.625,-0.084) (0.650,-0.079) (0.675,-0.073) (0.700,-0.067) (0.713,-0.063) (0.725,-0.060) (0.738,-0.056) (0.750,-0.051) (0.763,-0.046) (0.775,-0.041) (0.788,-0.035) (0.800,-0.027) (0.812,-0.019) (0.825,-0.009) (0.838,0.003) (0.850,0.016) (0.863,0.033) (0.875,0.052) (0.888,0.073) (0.900,0.096) (0.913,0.121) (0.925,0.145) (0.938,0.168) (0.950,0.189)};
% complete formula (center + distribution config), blue
\draw[very thick,blue] plot[smooth] coordinates {(0.050,-0.374) (0.062,-0.353) (0.075,-0.335) (0.088,-0.320) (0.100,-0.307) (0.113,-0.295) (0.125,-0.284) (0.138,-0.274) (0.150,-0.265) (0.162,-0.256) (0.175,-0.247) (0.188,-0.239) (0.200,-0.232) (0.213,-0.224) (0.225,-0.217) (0.237,-0.211) (0.250,-0.204) (0.275,-0.191) (0.300,-0.179) (0.325,-0.167) (0.350,-0.156) (0.375,-0.144) (0.400,-0.133) (0.425,-0.122) (0.450,-0.111) (0.475,-0.100) (0.500,-0.089) (0.525,-0.078) (0.550,-0.066) (0.575,-0.055) (0.600,-0.043) (0.625,-0.030) (0.650,-0.017) (0.675,-0.003) (0.700,0.011) (0.713,0.019) (0.725,0.027) (0.738,0.036) (0.750,0.044) (0.763,0.054) (0.775,0.064) (0.788,0.074) (0.800,0.085) (0.812,0.097) (0.825,0.110) (0.838,0.124) (0.850,0.140) (0.863,0.157) (0.875,0.175) (0.888,0.196) (0.900,0.218) (0.913,0.242) (0.925,0.268) (0.938,0.296) (0.950,0.328)};
% tab:qrev anchor points
\foreach \px/\py in {0.10/-0.307,0.25/-0.204,0.50/-0.089,0.75/0.044,0.90/0.218}
  {\fill (\px,\py) circle (1.3pt);}
% legend (top-left, clear of curves)
\draw[very thick,blue] (0.03,0.395) -- (0.115,0.395);
\node[font=\tiny,anchor=west] at (0.125,0.395) {완전식~\eqref{eq:complete} (중심$+$config)};
\draw[thick,densely dashed] (0.03,0.335) -- (0.115,0.335);
\node[font=\tiny,anchor=west] at (0.125,0.335) {단순식~\eqref{eq:weighted} (중심값만)};
\fill (0.073,0.275) circle (1.3pt);
\node[font=\tiny,anchor=west] at (0.125,0.275) {표~\ref{tab:qrev} 의 5점};
\node[font=\scriptsize] at (0.5,-0.56) {(a)};
\end{scope}
% =================== (b) Qrev/I(xbar) ===================
\begin{scope}[xshift=8.35cm,x=6.9cm,y=0.0182cm]
\draw[-{Latex[length=1.6mm]}] (-0.02,-112) -- (-0.02,126)
  node[above,font=\scriptsize] {$\dot Q_\rev/I=-T\,\partial U_\oc/\partial T$ [mV]};
\draw[-{Latex[length=1.6mm]}] (-0.02,-112) -- (1.06,-112)
  node[below,font=\scriptsize] {$\bar x$};
\foreach \xx in {0.2,0.4,0.6,0.8,1.0}
  {\draw (\xx,-110.2) -- (\xx,-113.8) node[below,font=\scriptsize] {\xx};}
\foreach \yy in {-100,-50,0,50,100}
  {\draw (-0.013,\yy) -- (-0.027,\yy) node[left,font=\scriptsize] {\yy};}
\draw[densely dashed,black!45] (-0.02,0) -- (1.04,0);
\draw[densely dotted,black!55] (0.681,-112) -- (0.681,30);
\node[font=\tiny,black!55,anchor=west] at (0.60,42) {$\bar x\!\approx\!0.68$};
% curve
\draw[very thick,blue] plot[smooth] coordinates {(0.050,111.42) (0.075,99.98) (0.100,91.53) (0.125,84.70) (0.150,78.88) (0.175,73.75) (0.200,69.10) (0.225,64.82) (0.250,60.81) (0.275,57.00) (0.300,53.36) (0.325,49.84) (0.350,46.42) (0.375,43.06) (0.400,39.74) (0.425,36.45) (0.450,33.17) (0.475,29.88) (0.500,26.57) (0.525,23.21) (0.550,19.80) (0.575,16.30) (0.600,12.70) (0.625,8.96) (0.650,5.06) (0.675,0.95) (0.700,-3.41) (0.725,-8.12) (0.750,-13.24) (0.775,-18.94) (0.800,-25.38) (0.812,-28.96) (0.825,-32.83) (0.838,-37.04) (0.850,-41.65) (0.863,-46.70) (0.875,-52.24) (0.888,-58.31) (0.900,-64.93) (0.913,-72.12) (0.925,-79.90) (0.938,-88.35) (0.950,-97.72)};
\foreach \px/\py in {0.10/91.5,0.25/60.8,0.50/26.6,0.75/-13.2,0.90/-64.9}
  {\fill (\px,\py) circle (1.3pt);}
% region labels
\node[font=\scriptsize,anchor=west] at (0.30,95) {방전($I{>}0$) \emph{발열} ($\Delta S<0$)};
\node[font=\scriptsize,anchor=west] at (0.42,-72) {\emph{흡열} ($\Delta S>0$, config 발산 지배)};
\node[font=\scriptsize] at (0.5,-134) {(b)};
\end{scope}
\end{tikzpicture}
\caption{완전식이 만드는 하프셀 엔트로피 계수와 가역 발열의 실계산 곡선 --- 좌표는 식 그대로의
수치 평가다(4-전이 staging 입력: 표~\ref{tab:ds} 의 $\Delta S^0_j$, $w_j{=}RT/F$($n_j{=}1$),
$298.15$ K; 각 $\bar x$ 에서 전하 보존 음함수~\eqref{eq:implicit} 를 풀어 $U_\oc$ 를 얻은 뒤 평가).
(a) $\partial U_\oc/\partial T(\bar x)$ --- 파란 실선 $=$ 완전식~\eqref{eq:complete}(중심 표준값$+$분포
config), 검은 파선 $=$ 단순식~\eqref{eq:weighted}(중심값만)이고 둘의 차가 봉우리 내부 config
항이다(파생 B). 점선 수평선은 전이별 수준 $\Delta S^0_j/F$ --- 곡선이 전이 경계마다 계단 없이 인접
수준 사이를 연속으로 지난다(파생 A 의 연속 블렌드). ● 는 표~\ref{tab:qrev} 의 다섯 점과 일치하며
$\bar x{=}0.25$ 값이 계산 예제(\S\ref{ssec:worked})의 $-0.204$($/$단순식 $-0.134$) mV/K 다.
(b) 가역 발열 $\dot Q_\rev/I=-T\,\partial U_\oc/\partial T$(식~\eqref{eq:qrev}) --- $\bar
x\!\approx\!0.68$ 에서 방전 발열이 흡열로 스스로 반전한다(임의 함수$\cdot$외부 스위치 없이 분포만으로).
완전식 곡선은 폭의 열적 서식 $w_j{=}n_jRT/F$ 를 전제하며, 실측 폭이 $T$-동결에 가까우면 단순식이
보수적 기준이다(\S\ref{ssec:weff}).}
\label{fig:cand-ff2-1}
\end{figure}
